\section{A deformation of the crossing identity on a Fermat curve}\label{sec:deformed-crossing}
The obstruction in Proposition~\ref{prop:fermion-grading-rigidity} fixes the finite fermion kernel. A different construction is possible: the binary factors in Lemma~\ref{lem:binary-crossing} admit sparse digit weights. Their vacuum normalizer becomes algebraic. We give the resulting exact identity on a curve, including its two local rank estimates. The global reconstruction cost is a separate issue, addressed in the following section. No improved crossing exponent follows from the local estimates alone.

The determinant construction and its auxiliary-block expansion are adaptations of \cite[Appendix B.3]{Anthropic2026}. The deformation below is proved directly; no priority claim is made.

Fix $q=2^p$ with $p\geq1$, and put $N=q-1$. For a nonnegative integer $i$, define
\[
 E_q(i)=\sum_{j\geq0}\operatorname{bit}_j(i)q^j,
 \qquad \ell_q(i)=2N E_q(i)+1.
\]
The operation $i\mathbin{\&}j$ denotes the integer whose binary digits are the products of the corresponding digits of $i$ and $j$. A symbol is a pair $(I,J)$ of finite subsets of nonnegative integers with $|I|=|J|$. Its energy is the sum of $\ell_q$ over both sets. Let $P_q(k)$ count symbols of energy at most $k$.

Work over $\mathbb F_2$ on the projective curve
\[
 \mathcal F_q:\quad X^N+Y^N=Z^N.
\]
For $N=1$ this is a line. For all $N$ it is smooth: its three partial derivatives cannot vanish simultaneously at a projective point, since $N$ is odd. It is absolutely irreducible. Indeed, in the affine equation $Y^N+(1+X^N)$, the prime $X-1$ divides the constant term exactly once; Eisenstein's criterion over the algebraic closure proves irreducibility. In the affine chart $Z=1$ set
\[
 y=X^{2N},\quad f=Y^2,\quad D_i=X^{\ell_q(i)},\quad
 B_{ij}=y^{E_q(j)-E_q(i\mathbin{\&}j)},\quad
 C_{ij}=y^{E_q(i\mathbin{\&}j)}.
\]
All these functions are regular on the affine curve. The points used below are
$P_0=[0:1:1]$, $P_\infty=[1:1:0]$, and $P_*=[1:0:1]$.

\begin{lemma}[Digit orthogonality on the distinguished branch]\label{lem:curve-orthogonality}
The formal series
\[
 \mathcal T_q(y)=\prod_{j\geq0}(1+y^{q^j})
\]
satisfies $\mathcal T_q(y)^{q-1}=(1+y)^{-1}$. At $P_0$ the coordinate $X$ is a uniformizer, $f=\mathcal T_q(y)^{-1}$, and
\begin{equation}\label{eq:curve-orthogonality}
 f\sum_{r\geq0}D_r B_{ri}B_{rj}=\delta_{ij}D_i.
\end{equation}
\end{lemma}
\begin{proof}
The product is a well-defined formal power series. Frobenius gives
$\mathcal T_q(y)=(1+y)\mathcal T_q(y^q)=(1+y)\mathcal T_q(y)^q$,
which proves the first assertion. The derivative of the curve equation with respect to $Y$ is one at $P_0$, so $Y$ has a unique power-series expansion with constant term one. Its square satisfies
$f^N=1+X^{2N}=1+y$. Since $N$ is odd, the $N$th root with constant term one is unique. This identifies $f$ with $\mathcal T_q^{-1}$.

For one binary digit, direct multiplication gives
\[
 \begin{pmatrix}1&1\\z&1\end{pmatrix}
 \begin{pmatrix}1&0\\0&z\end{pmatrix}
 \begin{pmatrix}1&z\\1&1\end{pmatrix}
 =(1+z)\begin{pmatrix}1&0\\0&z\end{pmatrix}.
\]
Multiply these identities over the digits, with $z=y^{q^j}$. Digits beyond those occurring in $i,j$ contribute the scalar $1+y^{q^j}$; a differing digit in $i,j$ contributes zero. Thus
$\sum_r y^{E_q(r)}B_{ri}B_{rj}=\mathcal T_q(y)\delta_{ij}y^{E_q(i)}$.
Multiplication by $Xf$ proves \eqref{eq:curve-orthogonality}. All infinite sums converge formally because $E_q(r)$ tends to infinity, while every $B_{ri}$ is a polynomial with nonnegative exponents.
\end{proof}

Fix the full nonempty classes $\Sigma_1,\Sigma_2$ of symbols of sizes $n_1,n_2$ and energies $e_1,e_2$. Full classes are used so that exchanging the two sets of a symbol permutes its class. Take labels $s_a,s_c\in\Sigma_1$ and $s_b,s_d\in\Sigma_2$, and write $s_v=(I_v,J_v)$. Define a square matrix $M$, with rows indexed by $I_a,I_d,I_c,I_b$ and columns by $J_a,J_d,J_c,J_b$, by the following block entries. Within each entry $i$ is the row index and $j$ the column index:
\[
\begin{array}{c|cccc}
 &a&d&c&b\\\hline
 a&0&C_{ij}&\delta_{ij}D_i&fD_iB_{ij}\\
 d&C_{ji}&0&D_jB_{ji}&\delta_{ij}D_i\\
 c&\delta_{ij}D_i&D_iB_{ij}&0&fC_{ij}\\
 b&fD_jB_{ji}&\delta_{ij}D_i&fC_{ji}&0
\end{array}
\]
Let $\Psi=\det M$, a tensor-valued regular function, and put $\Theta=2(e_1+e_2)$.

\begin{theorem}[Crossing, reflection, and the two local rank bounds]\label{thm:curve-crossing}
The value $\Psi(P_*)$ is the product of the two crossing wires of dimensions $|\Sigma_1|,|\Sigma_2|$, after exchanging $I,J$ in the labels at $c,d$. The involution
\[
 \sigma(X,Y)=(1/X,Y/X)
\]
interchanges $P_0,P_\infty$ and satisfies
\begin{equation}\label{eq:curve-reflection}
 X^\Theta\Psi(\sigma(X,Y);s_a,s_b,s_c,s_d)
 =\Psi(X,Y;s_a,s_d,s_c,s_b).
\end{equation}
The coefficient of $X^k$ in the expansion at $P_0$ has matrix rank at most $(k+1)P_q(k)$ across $\{a,d\}\mid\{b,c\}$. The coefficient of $t^k$ in the expansion of $X^{-\Theta}\Psi$ at $P_\infty$, using $t=1/X$, has the same bound across $\{a,b\}\mid\{c,d\}$.
\end{theorem}
\begin{proof}
At $P_*$, $f=0$, $D_i=1$ and $B_{ij}=C_{ij}=1$. Partition the rows and columns into the pairs $\{a,d\}$ and $\{c,b\}$. The second diagonal block vanishes. Each block between the pairs is triangular with diagonal blocks given by the appropriate Kronecker deltas. Its determinant is consequently
\[
 [I_a=J_c][I_c=J_a][I_b=J_d][I_d=J_b],
\]
as asserted.

The projective transformation $[X:Y:Z]\mapsto[Z:Y:X]$ preserves the curve and restricts to $\sigma$. It sends $y$ to $1/y$ and $f$ to $f/X^2$. Moreover
\[
 B_{ij}\circ\sigma=y^{-E_q(j)}C_{ij},\qquad
 C_{ij}\circ\sigma=y^{-E_q(j)}B_{ij},\qquad
 D_i\circ\sigma=D_i^{-1}.
\]
In $M\circ\sigma$, multiply both the row and column at index $i$ of parties $a,c$ by $D_i$, of party $d$ by $f y^{E_q(i)}$, and of party $b$ by $X^2y^{E_q(i)}$. Substitution in the displayed block matrix gives $M$ with $b,d$ exchanged, and with both rows and columns of the old party $b$ additionally multiplied by $f$. The total multiplier on the left determinant is
$X^{2e_1+(e_2-2n_2)+(e_2+2n_2)}f^{2n_2}=X^\Theta f^{2n_2}$.
The extra multiplier on the right is $f^{2n_2}$. Cancelling this nonzero function proves \eqref{eq:curve-reflection}.

For the coefficient bound at $P_0$, keep only auxiliary indices $r$ with $\ell_q(r)\leq k$, and let $D$ be the corresponding diagonal matrix with entries $D_r$. Let $F$ have rows $(\delta_{ir})_r$ at $a$ and $(B_{ri})_r$ at $d$. Let $G$ have columns $(\delta_{rj})_r$ at $c$ and $(fB_{rj})_r$ at $b$. Orthogonality shows that the block between these two pairs is congruent to $FDG$ modulo $X^{k+1}$. The reverse block is congruent to $G'DF'$, with the corresponding transposed choices. All omitted terms have order greater than $k$.

Write the two diagonal blocks as $M_1,M_2$. The block determinant identity replaces
$\left(\begin{smallmatrix}M_1&FDG\\G'DF'&M_2\end{smallmatrix}\right)$,
without division by any power of $X$, by the auxiliary determinant
\[
 \det\begin{pmatrix}
 M_1&0&F&0\\
 0&M_2&0&G'\\
 0&DG&1&0\\
 DF'&0&0&1
 \end{pmatrix}.
\]
This is the Schur-complement identity over the power-series ring; signs disappear in characteristic two. Expand each of its last block rows into its $D$-part and its identity part. Choosing the $D$-parts on subsets $T_1,T_2$ contributes the factor
$X^{\sum_{T_1}\ell_q+\sum_{T_2}\ell_q}$.
After the unused identity rows and columns are removed, the remaining determinant is a product of a function of $(s_a,s_d)$ and a function of $(s_b,s_c)$. It vanishes unless $|T_1|=|T_2|$, since otherwise the two separated blocks are rectangular. Thus at most $P_q(k)$ symbols can contribute to the coefficient of degree $k$. For each such symbol, splitting the remaining coefficient degree between the two factors uses at most $k+1$ products. This proves the first rank bound. Applying \eqref{eq:curve-reflection} proves the second one.
\end{proof}

\begin{corollary}[The full pole divisor]\label{cor:curve-poles}
There are $N$ distinct points at infinity over $\mathbb F_q$. The function $\Psi$ has pole order at most $\Theta$ at every one of them and has no affine poles.
\end{corollary}
\begin{proof}
The infinity points are $[1:\zeta:0]$, with $\zeta^N=1$, and $X$ has a simple pole at each. Under $\sigma$ they are the points $[0:\zeta:1]$. The determinant is regular at every affine point. Formula \eqref{eq:curve-reflection} therefore gives the asserted pole bound at all infinity points, not just at $P_\infty$.
\end{proof}

The local orthogonality proof singles out $P_0$. At another zero of $X$, with $Y=\zeta\ne1$, the value of the normalizer is multiplied by $\zeta^2$. The auxiliary block product then has a residual $(1+\zeta^2)D$ term. Consequently the same low-rank statement at every other branch has not been established. Ignoring these other poles would omit the global reconstruction problem.

\begin{lemma}[Exact interleaving for the sparse energies]\label{lem:curve-interleaving}
For nonnegative integer bounds $A,B$,
\[
 P_q(q(A+B))\geq P_q(A)P_q(B).
\]
\end{lemma}
\begin{proof}
The digit definition gives
$\ell_q(2r)=q\ell_q(r)-N$ and $\ell_q(2r+1)=q\ell_q(r)+N$.
Given symbols $(I_1,J_1)$ and $(I_2,J_2)$, put
\[
 I=2I_1\,\cup\,(2I_2+1),\qquad
 J=(2J_1+1)\,\cup\,2J_2.
\]
This is an injective map to symbols. Within each input symbol the positive and negative offsets cancel because its two sets have equal cardinality. Its output energy is exactly $q$ times the sum of the input energies.
\end{proof}

This interleaving does not by itself prove a limit of the form $\log P_q(E)\sim\kappa E^\alpha$. Sparse digit energies can produce scale-dependent leading factors. The finite bounds in Theorem~\ref{thm:curve-crossing} use the exact counting function and require no such asymptotic claim.

\subsection{A counting bound without a smooth entropy asymptotic}
The sparse energies need not have a constant leading entropy factor. Concavity of a suitable envelope suffices to test the generic counting budget. Natural logarithms are used in this subsection.

\begin{theorem}[Concave envelope for the symbol count]\label{thm:curve-count-envelope}
Fix $q=2^p$ and put $\alpha=1/(p+1)$. There is an increasing concave function $S_q:[0,\infty)\to[0,\infty)$, with $S_q(0)=0$, such that
\[
 S_q(E)=\Theta(E^\alpha),\qquad
 \ln P_q(E)=S_q(E)+o(E^\alpha)\quad(E\longrightarrow\infty).
\]
Consequently, for every fixed $0<\lambda\leq1$,
\[
 \ln P_q(\lambda E)\geq(\lambda-o(1))\ln P_q(E),
\]
and
\begin{equation}\label{eq:curve-count-ratio}
 \liminf_{E\to\infty}
 \frac{\ln P_q(2(q-1)E)}{\ln P_q(E)}\geq2-\frac2q.
\end{equation}
\end{theorem}
\begin{proof}
First count single finite subsets of the mode indices. Their weights satisfy
$\ell_q(r)\asymp(r+1)^p$,
because $2^j\leq r<2^{j+1}$ implies
$q^j\leq E_q(r)\leq(q^{j+1}-1)/(q-1)$; the index zero is absorbed into the constants. Let $\beta=1/p$ and define, for $t>0$,
\[
 H(t)=\sum_{r\geq0}\ln(1+e^{-t\ell_q(r)}),\qquad
 m(t)=\sum_{r\geq0}\frac{\ell_q(r)}{1+e^{t\ell_q(r)}}.
\]
Comparison with the integrals of $e^{-ct x^p}$, $x^pe^{-ct x^p}$ and $x^{2p}e^{-ct x^p}$ gives
\[
 H(t)=\Theta(t^{-\beta}),\qquad
 m(t)=\Theta(t^{-\beta-1}),\qquad
 v(t):=\sum_r\frac{\ell_q(r)^2e^{t\ell_q(r)}}{(1+e^{t\ell_q(r)})^2}
 =O(t^{-\beta-2})
\]
as $t\downarrow0$. For the lower bounds, restrict to a fixed interval of $\Theta(t^{-\beta})$ indices for which the weights lie between two positive constant multiples of $1/t$. For the upper bounds, use the displayed exponential comparisons. The sums converge locally uniformly for $t>0$, so differentiation is permitted. In particular $m=-H'$ is continuous and strictly decreasing from infinity to zero.

Set
\[
 S_1(E)=\inf_{t>0}\{tE+H(t)\}.
\]
This is increasing and concave, with $S_1(0)=0$. Its minimizing parameter is determined by $m(t)=E$, so $t\asymp E^{\alpha-1}$ and $S_1(E)=\Theta(E^\alpha)$.

Let $A(E)$ count finite subsets of total weight at most $E$. The generating product gives $\ln A(E)\leq S_1(E)$. For the matching estimate, use independent Bernoulli occupancies with probabilities $(1+e^{t\ell_q(r)})^{-1}$. They occupy only finitely many indices almost surely, since the sum of these probabilities converges. A subset of weight $w$ has probability $e^{-tw-H(t)}$; the mean and variance of its weight are $m(t),v(t)$.

Choose $\delta=E^{-\alpha/4}$ and the parameter with $m(t)=(1-\delta)E$. Then $v(t)=O(E^{2-\alpha})$. Chebyshev's inequality shows that the weight lies within $\delta E/2$ of its mean with probability $1-O(E^{-\alpha/2})$. This interval lies below $E$. Counting its states gives
\[
 \ln A(E)\geq S_1((1-\delta)E)-t\delta E/2+o(1).
\]
The variational definition also gives
$S_1(E)\leq S_1((1-\delta)E)+t\delta E$.
Since $t\delta E=O(E^{3\alpha/4})=o(E^\alpha)$, we obtain
$\ln A(E)=S_1(E)+o(E^\alpha)$.

Now impose equal cardinality on a pair of subsets. The unrestricted two-colour generating function is $\exp(2H(t))$, giving
\[
 \ln P_q(E)\leq\inf_t\{tE+2H(t)\}=2S_1(E/2).
\]
If $A_j(E/2)$ counts subsets with cardinality $j$ and weight at most $E/2$, then
\[
 P_q(E)\geq\sum_j A_j(E/2)^2
 \geq\frac{A(E/2)^2}{\lfloor E/2\rfloor+1}.
\]
The last denominator suffices because every mode has weight at least one. Hence
$\ln P_q(E)=2S_1(E/2)+o(E^\alpha)$.
Take $S_q(E)=2S_1(E/2)$. Its concavity and its value at zero imply
$S_q(\lambda E)\geq\lambda S_q(E)$ for $0<\lambda\leq1$, proving the first consequence.

Finally put $\lambda=(q-1)/q$ and $a=\lfloor\lambda E\rfloor$. Lemma~\ref{lem:curve-interleaving} gives
$P_q(2(q-1)E)\geq P_q(2qa)\geq P_q(a)^2$.
Concavity gives $S_q(a)\geq(a/E)S_q(E)$, and $a/E=\lambda+o(1)$. Dividing by $\ln P_q(E)=\Theta(E^\alpha)$ proves \eqref{eq:curve-count-ratio}.
\end{proof}

For every fixed radix $q\geq4$, the ratio in \eqref{eq:curve-count-ratio} is at least $3/2>\sqrt2$. Thus a reconstruction certificate charged by a polynomial factor times $P_q(2(q-1)E)$ cannot, on the strength of that assigned budget, certify a crossing exponent below $\sqrt2$. A fixed-energy class has cardinality at most $P_q(E)$, which only strengthens this comparison. This is a limitation of the generic count estimate; it is not a lower bound for actual coefficient ranks or for the smaller span of the particular determinants. The constants here depend on the fixed radix, and no uniform assertion for a radix growing with $E$ is made.

\subsection{What the Galois trace retains}
Taking a trace gives a polynomial in $X$, but may discard the local low-rank property that made the curve construction useful.

\begin{proposition}[Trace collapse in the first symbol sizes]\label{prop:curve-trace-collapse}
Replace $f$ in the matrix by an independent variable $u$, and write the resulting determinant as $\Psi(X,u)$. Then $\deg_u\Psi\leq2\min(n_1,n_2)$, and its constant term is $X^{e_1+e_2}$ times the crossing tensor. On the curve, trace over $Y\mapsto\zeta Y$, $\zeta^N=1$, preserves the crossing evaluation and yields a polynomial in $X$ of degree at most $\Theta$. If $2\min(n_1,n_2)<N$, this trace is exactly its constant term in $u$. Its only nonzero coefficient has adjacent flattening rank $|\Sigma_1||\Sigma_2|$ across both cuts.
\end{proposition}
\begin{proof}
Consider a nonzero permutation monomial in the determinant. Let $\alpha$ count entries between $a,c$ in either direction, and $\beta$ entries between $b,d$. Counting rows and columns in the pair $\{a,c\}$ gives
$2n_1-\alpha=2n_2-\beta$.
Every entry incident to $b$ contributes $u$ except those between $b,d$. Thus its $u$-degree is
$h=2n_2-\beta=2n_1-\alpha\leq2\min(n_1,n_2)$.
When $h=0$, only the two crossing-wire blocks can occur. Their Kronecker deltas give the crossing tensor and their powers of $X$ sum to $e_1+e_2$.

The trace of $u^h=f^h=Y^{2h}$ is zero unless $N$ divides $h$, in which case it equals $f^h$: the sum of $\zeta^{2h}$ over the $N$ roots of unity is zero or one, respectively. The trace therefore lies in $\mathbb F_2[X]$, using $f^N=1+X^{2N}$. The pole bound in Corollary~\ref{cor:curve-poles} bounds its degree by $\Theta$. At $P_*$ every conjugate has the same crossing value, and the odd number $N$ equals one in the base field. If the degree in $u$ is smaller than $N$, only $h=0$ survives. Each of the two crossing wires crosses each adjacent partition, so both flattening ranks equal the product of their dimensions.
\end{proof}

The first trace beyond this collapse condition was also computed exactly. The program uses arithmetic in
$\mathbb F_2[X,f]/(f^N+1+X^{2N})$, expands determinants by a subset recurrence, and extracts the trace by retaining the reduced $f^0$ coefficient. It verifies the crossing evaluation entry by entry. The linear case is a control against the earlier calculation.
\begin{center}
\begin{tabular}{rrrrrr}
radix & size & $E$ & $Q$ & terms & max. smaller rank\\\hline
$2$ & $1$ & $4$ & $2$ & $8$ & $2$\\
$4$ & $2$ & $64$ & $8$ & $65$ & $64$
\end{tabular}
\end{center}
Here $E$ is the symbol energy, $Q$ the wire dimension, and the last column is the maximum over coefficients of the smaller of the two adjacent flattening ranks. The second row uses all eight labels whose four digit energies sum to ten, with available mode digit energies $0,1,4,5$. Its coefficient of degree $128$ is exactly the original crossing tensor, entry by entry; both adjacent flattening ranks there are $64$. Its bridge ratio is $\log_8 64=2$, so this trace supplies no strict shortening through the two-hub criterion. These are exact finite examples, not an exhaustive test of symbol classes or a lower bound for all deformed identities. The programs, complete rank profiles, and input labels are included in the workbench certificates.

\subsection{An exact expansion at the other branches}
The loss of orthogonality at the other branches can be expressed by finitely many deleted-row and deleted-column determinants. This isolates the algebraic terms that a recursive rank estimate would have to control.

\begin{proposition}[Bad-branch defect expansion]\label{prop:curve-defects}
Let $M_z$ be the displayed crossing matrix with $f$ replaced by $zf$, where $z^N=1$, and let $M_1$ denote the matrix on the distinguished branch. For
\[
 S\subseteq I_d\cap J_b,\qquad T\subseteq I_b\cap J_d,
\]
put $d(S,T)=|S|+|T|$ and
$w(S,T)=\sum_{i\in S}\ell_q(i)+\sum_{i\in T}\ell_q(i)$.
Let $M_1[S,T]$ be obtained by deleting rows $(d,S),(b,T)$ and columns $(b,S),(d,T)$. Then
\begin{equation}\label{eq:curve-defect-expansion}
 \det M_z=\sum_{S,T}
 z^{2n_2-d(S,T)}(1+z)^{d(S,T)}
 X^{w(S,T)}\det M_1[S,T].
\end{equation}
Consequently its root-of-unity trace is
\begin{equation}\label{eq:curve-defect-trace}
 \sum_{z^N=1}\det M_z
 =\sum_{S,T}a_{d(S,T)}X^{w(S,T)}\det M_1[S,T],
 \qquad
 a_d=\sum_{j=0}^d\binom dj[\,N\mid2n_2-d+j\,]
 \quad\text{in }\mathbb F_2.
\end{equation}
The residual labels at $b,d$ both have energy $e_2-w(S,T)$. Their charges, defined as $|I|-|J|$, are respectively $|S|-|T|$ and $|T|-|S|$. Thus the minor has ordinary balanced symbol labels at these parties precisely when $|S|=|T|$.
\end{proposition}
\begin{proof}
Let $G_z$ be diagonal, with entry $z$ on every row and column belonging to $b$, and entry one elsewhere. Let $R_{bd}$ have entries $\delta_{ij}D_i$ in the $b,d$ and $d,b$ blocks and zero elsewhere. Checking the four blocks incident with $b$ gives
\[
 M_z=G_zM_1G_z+(1+z)R_{bd}.
\]
In the determinant expansion, select the entries drawn from the second summand. They form a partial matching: a subset $S$ uses rows of $d$ and matching columns of $b$, while a subset $T$ uses rows of $b$ and matching columns of $d$. Such entries are available exactly for the intersections stated above. Their product is $(1+z)^{|S|+|T|}X^{w(S,T)}$. The remaining minor of $G_zM_1G_z$ has $n_2-|T|$ remaining $b$ rows and $n_2-|S|$ remaining $b$ columns, hence the scalar factor $z^{2n_2-|S|-|T|}$. Signs vanish in characteristic two. Summing over all selected partial matchings proves \eqref{eq:curve-defect-expansion}.

For any integer $h$, the sum of $z^h$ over the $N$th roots of unity is zero unless $N\mid h$, and equals $N=1$ in the field when $N\mid h$. Expanding $(1+z)^d$ therefore gives \eqref{eq:curve-defect-trace}. At $b$ the surviving row and column sets are $I_b\setminus T,J_b\setminus S$; at $d$ they are $I_d\setminus S,J_d\setminus T$. This proves both the energy and the charge formulas.
\end{proof}

For $N=3,n_2=2$, the trace weights $(a_0,a_1,a_2,a_3,a_4)$ are $(0,1,0,1,1)$. Every surviving proper partial deletion therefore has total size one or three, with opposite charges $+1,-1$ on $b,d$. The only possible surviving balanced nonempty deletion has total size four, which removes all rows and columns of those parties; it occurs only when the required intersections contain both indices. This finite identity does not close a recurrence on the original balanced symbol classes. A useful recursion would also need rank bounds for the charged minors and a quantitative count of the resulting terms and their selector costs.

\begin{proposition}[The coarse charged-minor bound]\label{prop:charged-minors}
Let $P_{q,\delta}(k)$ count pairs $(U,V)$ of finite mode sets with $|U|-|V|=\delta$ and total energy at most $k$. For fixed defects $S,T$, the coefficient of $X^k$ in $\det M_1[S,T]$ has rank at most
\[
 (k+1)P_{q,|T|-|S|}(k)
\]
across $ad\mid bc$. Let $Q_q(k)$ count quadruples $(S,T,U,V)$ with total energy at most $k$ and $|S|+|U|=|T|+|V|$. At any branch $z^N=1$, the corresponding coefficient of $\det M_z$ has rank at most $(k+1)Q_q(k)$. For fixed $q=2^p$, with $\alpha=1/(p+1)$ and the envelope from Theorem~\ref{thm:curve-count-envelope},
\[
 \ln Q_q(E)=2S_q(E/2)+o(E^\alpha),\qquad
 \lim_{E\to\infty}\frac{\ln Q_q(2E)}{\ln P_q(E)}=2.
\]
\end{proposition}
\begin{proof}
In the residual matrix, the $ad$ block has $r_1=n_1+n_2-|S|$ rows and $c_1=n_1+n_2-|T|$ columns. Apply the auxiliary determinant expansion from Theorem~\ref{thm:curve-crossing}; it does not require the two diagonal blocks to be square. If its third and fourth block rows select the diagonal factors on $U,V$, respectively, the separated first block has $r_1+|V|$ rows and $c_1+|U|$ columns. It can be nonzero only when
$|U|-|V|=r_1-c_1=|T|-|S|$.
The auxiliary monomial has the total energy of $U,V$, and all remaining entries have nonnegative valuation. Splitting the remaining degree uses at most $k+1$ products of functions on the two sides. This proves the first assertion. Fixed deletions are local linear relabelings of the terminal spaces and cannot increase this flattening rank.

Sum the rank bounds over the defects in \eqref{eq:curve-defect-expansion}. Their energy plus the auxiliary energy is at most $k$, and the cardinality condition is precisely $|S|+|U|=|T|+|V|$. Discarding the restrictions imposed by the original symbol classes gives the asserted coarse bound using $Q_q$.

For the asymptotic count, two independently balanced symbols, each of energy at most $E/2$, give a quadruple counted by $Q_q(E)$. Hence
$Q_q(E)\geq P_q(E/2)^2$.
For the upper bound discard the cardinality condition and use the one-colour function $H$ from the proof of Theorem~\ref{thm:curve-count-envelope}:
\[
 Q_q(E)\leq\inf_{t>0}\exp(tE+4H(t))
 =\exp(4S_1(E/4))=\exp(2S_q(E/2)).
\]
The already proved asymptotic for $P_q$ supplies the matching lower logarithm. Substituting $2E$ and dividing by $\ln P_q(E)=\Theta(E^\alpha)$ proves the limit.
\end{proof}

Thus this direct treatment of all defects assigns an exponent of two at the budget $2E$. It gives no improvement over the full rank of the crossing. This is a limitation of the coarse count, not a lower bound on the actual coefficient ranks: the fixed symbol classes, the trace weights, and cancellations between defects have been discarded in its derivation.

\begin{proposition}[The highest normalizer coefficient]\label{prop:top-normalizer-coefficient}
Suppose both symbol classes have size $n$, and retain the normalizer as an independent variable:
\[
 \Psi(X,u)=\sum_{h=0}^{2n}u^h A_h(X).
\]
At the distinguished branch, with $f(0)=1$, the top coefficient satisfies
\begin{equation}\label{eq:top-normalizer-defects}
 f^{2n}A_{2n}(X)
 =\sum_{S,T}X^{w(S,T)}\det M_1[S,T].
\end{equation}
Here the defects and their energies are those of Proposition~\ref{prop:curve-defects}. Consequently, for every $k\geq0$,
\[
 \operatorname{rank}_{ad\mid bc}[X^k]A_{2n}
 \leq(k+1)^2Q_q(k).
\]
If $G\in\mathbb F_2[X]$ divides every coordinate of $A_{2n}$ and $G(0)=1$, the corresponding bound for $A_{2n}/G$ is $(k+1)^3Q_q(k)$. If both input classes have energy $E$, then
\begin{equation}\label{eq:top-normalizer-reflection}
 X^{4E-4n}A_{2n}(1/X;s_a,s_b,s_c,s_d)
 =A_{2n}(X;s_a,s_d,s_c,s_b).
\end{equation}
\end{proposition}
\begin{proof}
The determinant expansion proving \eqref{eq:curve-defect-expansion} uses the relation $z^N=1$ only for its subsequent trace. It is therefore valid with an independent indeterminate $z$. Its scalar factor for defects of total size $d$ is $z^{2n-d}(1+z)^d$, whose coefficient of $z^{2n}$ is one. Extracting that coefficient from $\Psi(X,zf)$ proves \eqref{eq:top-normalizer-defects}.

The charged-minor proof bounds each degree-$j$ coefficient of the right side by $(j+1)Q_q(j)$. Since $f^{2n}$ is a scalar unit in the power-series ring at $X=0$, multiplication by its inverse expresses the degree-$k$ coefficient of $A_{2n}$ as a sum of at most $k+1$ such earlier coefficients. Rank subadditivity and monotonicity of $Q_q$ give $(k+1)^2Q_q(k)$. A second scalar-unit convolution with $G^{-1}$ gives the stated cubic factor.

Finally, the row-and-column proof of \eqref{eq:curve-reflection} is valid with $u$ independent, provided the involution sends $u$ to $u/X^2$. It gives
\[
 X^{4E}\Psi(1/X,u/X^2;s_a,s_b,s_c,s_d)
 =\Psi(X,u;s_a,s_d,s_c,s_b).
\]
Comparing coefficients of $u^{2n}$ proves \eqref{eq:top-normalizer-reflection}.
\end{proof}

Thus a common scalar divisor can reduce the degree needed for polynomial reconstruction while preserving an explicit coefficient bound. The bound above retains all defect choices; improving the actual ranks requires additional information about their cancellations. It does not assert that the normalized top coefficient has a full crossing-wire output.


\subsection{Character projections at the output point}
Character traces can select residue classes of powers of $f$. Their values at the crossing point are particularly restrictive.

\begin{proposition}[Only the trivial character retains this output]\label{prop:character-output}
Use fixed-size, equal-energy classes at opposite parties, with injective individual mode energies. Specialize $D_i=B_{ij}=C_{ij}=1$ but retain $f=u$ as an independent indeterminate. The determinant is the crossing tensor for every $u$ in characteristic two.

Write the unspecialized determinant as $\Psi(X,f)=\sum_h f^h A_h(X)$. For $0\leq r<N$, the divided character trace
\[
 T_r(X)=f^{-r}\sum_{z^N=1}z^{-r}\Psi(X,zf)
 =\sum_{h\equiv r\;({\rm mod}\,N)}
 (1+X^{2N})^{(h-r)/N}A_h(X)
\]
is a polynomial. At $X=1$, $T_0$ is the crossing tensor and $T_r=0$ for every $r>0$.
\end{proposition}
\begin{proof}
Group rows and columns by the opposite pairs $ac$ and $bd$. The two diagonal blocks, denoted $A,B$, contain only the Kronecker-delta matchings between opposite parties. The off-diagonal blocks have rank one: all entries on the $ac$ side are one, while the $bd$ side has weights one at $d$ and $u$ at $b$. The block determinant formula, valid as a polynomial identity, is
\[
 \det\begin{pmatrix}A&av^{\mathsf T}\\wa'^{\mathsf T}&B\end{pmatrix}
 =\det A\det B+
 (a'^{\mathsf T}\operatorname{adj}(A)a)
 (v^{\mathsf T}\operatorname{adj}(B)w),
\]
where $a,a'$ are the respective all-ones vectors. It follows, for example, by applying the rank-one determinant lemma when $A,B$ are invertible and then clearing denominators.

We show that the first adjugate contraction vanishes. If $A$ is invertible, it is a permutation matrix and the contraction is $2n_1=0$. If its rank is at most $2n_1-2$, its adjugate is zero. The remaining rank $2n_1-1$ cannot occur. Indeed, the two Kronecker blocks then have ranks $n_1$ and $n_1-1$. One opposite matching is exact, while the other two sets differ by one mode. Equality of the two symbol energies makes those two modes have equal weight, contradicting injectivity unless the sets agree. When $n_1=0$ the determinant assertion is immediate. Hence the update term is always zero, and $\det A\det B$ is exactly the crossing.

It follows that $A_0(1)$ is the crossing and $A_h(1)=0$ for $h>0$. The root-of-unity sum retains exactly $h\equiv r$; division by $f^r$ leaves a nonnegative power of $f^N=1+X^{2N}$. At $X=1$ only $h=r$ can survive, giving the stated values.
\end{proof}

This excludes using a single nontrivial character trace with the same output evaluation. It does not exclude a combination with the trivial character, a derivative at the output point, or a different linear reconstruction functional.


\paragraph{A higher local coefficient in the eight-label example.}
For radix four, size two and energy $64$, exact calculation gives the following character data. ``Order'' is the common order of vanishing at $X=1$; the primitive polynomial is obtained by dividing by the complete polynomial gcd of its coordinates.
\begin{center}
\begin{tabular}{@{}crrrrrc@{}}
\toprule
Char. & Span & Order & Gcd deg. & Degree & Wire & Ranks\\
\midrule
$0$&$13$&$0$&$38$&$180$&$8$&$0+64$\\
$1$&$28$&$42$&$62$&$192$&$8$&$1+28$\\
$2$&$16$&$4$&$68$&$120$&$2$&$0+4$\\
\bottomrule
\end{tabular}
\end{center}
The wire-rank column gives the rank of each of the two opposite-pair matrices in the primitive output at $X=1$. The grouped ranks are the best pair found among all contiguous low-degree/high-degree splits, using the two adjacent cuts; they are not an optimization over all decompositions.

For character one, the gcd is exactly
$(1+X^6)^5(1+X)^{32}$.
The primitive output is a product of two permutation wires, each with permutation $P=(2\ 5)(3\ 4)$ in the zero-based eight-label ordering of Proposition~\ref{prop:finite-two-bit}. Its constant coefficient is the other pairing, with permutation $Q=(3\ 4)$ on both wires. After normalizing the output wires, this parallel pairing acts across $ab\mid cd$ as
$(R\otimes R)\mathsf F$, where $R=PQ=(2\ 5)$ and $\mathsf F$ swaps the two factors. Since $R^2=1$, this operator is conjugate to $\mathsf F$. The remainder is therefore conjugate to $1+\mathsf F$, whose rank in characteristic two is $\binom82=28$: it vanishes on diagonal basis vectors and has rank one on each two-dimensional off-diagonal pair. Thus this particular recovery is the usual exterior-square routing identity. It provides no improved asymptotic crossing estimate. A nonzero polynomial $2\times2$ minor also shows that the character-one tensor has opposite-cut rank greater than one over $\mathbb F_2(X)$; the product of two wires occurs only in its leading local output.

\subsection{Finite truncations of the digit kernel}
A finite product makes the normalizer rational, but also changes the crossing specialization. The following elementary facts isolate this issue.

\begin{proposition}[Finite normalization and reciprocal scaling]\label{prop:finite-kernel-gates}
Fix $L\geq1$, let $M=2^L$, and put
$H=1+q+\cdots+q^{L-1}$. For $0\leq i<M$ use the same digit energies $E_q(i)$ and set
\[
 \mathcal T_L(y)=\prod_{j<L}(1+y^{q^j})=(1+y)^H,
 \qquad f_L(y)=\mathcal T_L(y)^{-1}.
\]
Then
\[
 f_L(y)\sum_{r<M}y^{E_q(r)}B_{ri}B_{rj}
 =\delta_{ij}y^{E_q(i)},\qquad
 f_L(1/y)=y^Hf_L(y).
\]
The normalizer has no finite zero. More generally, suppose a finite sum satisfies
\[
 f\sum_rD_rB_{ri}B_{rj}=\delta_{ij}D_i.
\]
There is no specialization at which $f=0$, all its entries are regular, and one target $D_i$ has nonzero value.

For the fixed crossing block pattern, let $y=X^{2s}$ and $D_i=X^\beta y^{E_q(i)}$. Suppose inversion $X\mapsto1/X$, followed by arbitrary nonzero row and column multipliers, gives the same block matrix with $b,d$ exchanged, without permuting mode indices. If all four parties contain mode zero, this requires
\[
 \beta=-sH.
\]
These centered weights have both positive and negative exponents.
\end{proposition}
\begin{proof}
Frobenius gives $1+y^{q^j}=(1+y)^{q^j}$, proving the product formula. The finite product of the binary Gram identities in Lemma~\ref{lem:curve-orthogonality} proves the finite sum. Reciprocity follows by inverting every factor. The obstruction to a regular specialization is immediate on setting $i=j$: its left side vanishes and its right side does not.

For the scaling assertion write $t=X^\beta$ and $\kappa=f_L(1/y)/f_L(y)=y^H$. At mode zero, all $B,C$ entries are one. The cross ratio using rows $a,b$ and columns $c,d$ is invariant under arbitrary row and column multiplication. For the inverted matrix it is
\[
 \frac{M_{ac}M_{bd}}{M_{ad}M_{bc}}
 =\frac1{\kappa f_Lt^2},
\]
whereas the matrix with $b,d$ exchanged has cross ratio $1/f_L$. Thus $\kappa t^2=1$, which is exactly $\beta=-sH$. Since $E_q(0)=0$ and $E_q(M-1)=H$, the exponents $2sE_q(i)-sH$ have both signs.
\end{proof}

The regular-specialization obstruction concerns the original output mechanism. For a particular finite family, the crossing can also fail to belong to the span of every coefficient tensor, as the next exact identity shows.

\begin{proposition}[A two-bit coefficient-span obstruction]\label{prop:finite-two-bit}
Use four modes $0,1,2,3$, with
$a_i=x^{\operatorname{bit}_1(i)}y^{\operatorname{bit}_0(i)}$,
$D_i=ta_i$, $C_{ij}=a_{i\mathbin{\&}j}$ and
$B_{ij}=a_j/a_{i\mathbin{\&}j}$. Let the eight symbol labels, in order, be
\[
\begin{array}{c|cccc}
 j&0&1&2&3\\\hline
 I_j&01&02&03&03\\
 J_j&23&13&03&12
\end{array}
\qquad
\begin{array}{c|cccc}
 j&4&5&6&7\\\hline
 I_j&12&12&13&23\\
 J_j&03&12&02&01
\end{array}
\]
where, for example, $01$ denotes $\{0,1\}$. Retain $f$ as an independent variable in the crossing matrix and write $\Psi_{abcd}$ for its determinant at these labels. Over $\mathbb F_2[f,t,x,y]$,
\begin{align}\label{eq:two-bit-annihilator}
 &\Psi_{7610}+\Psi_{7660}+\Psi_{7700}+\Psi_{7770}\nonumber\\
 &\quad=\bigl[1+f^2(1+x^2)(1+y^2)\bigr]\nonumber\\
 &\qquad\quad\cdot\bigl[t^4f^2x^4y^2(1+y^2)
       +t^8x^4y^4\{1+f^2x^2(1+y^2)\}\bigr].
\end{align}
For the finite normalizer $f=((1+x)(1+y))^{-1}$, the crossing tensor is therefore outside the span of all coefficient tensors and regular values of this rational family.
\end{proposition}
\begin{proof}
Expanding the four determinants, their sum is
\begin{align*}
 t^8x^4y^4
 &+f^2\bigl[t^4x^4(y^2+y^4)+t^8x^4(y^4+y^6)\bigr]\\
 &+f^4\bigl[t^4x^4(y^2+y^6)(1+x^2)
             +t^8x^6y^4(1+y^4)(1+x^2)\bigr].
\end{align*}
This expansion has thirteen monomials; multiplying the right side of \eqref{eq:two-bit-annihilator} gives the same expression. With the stated normalizer the first factor vanishes. Hence the sum of these four coordinates annihilates the complete rational family, and therefore every coefficient or regular value. On the crossing tensor these four coordinates are $0,0,1,0$, respectively. Their sum is one, proving the exclusion.
\end{proof}

The annihilator fixes the terminal coordinates. A locally equivalent product of two full-rank wires would also be a useful output, and is not excluded by this linear test. The displayed obstruction is specific to the four-mode family. Keeping the same eight labels and adding four unused auxiliary modes changes the normalizer and removes it. For example, with radix four and $D_i=X^{6E_4(i)+1}$, exact elimination gives coefficient-span dimensions $24$ and $45$ for four and eight auxiliary modes, respectively. The crossing is absent in the first span and present in the second. In the second case, after clearing the common denominator $(1+X^6)^{84}$, the coefficient of degree $512$ is exactly the crossing and has adjacent flattening ranks $64$. Thus this recovery does not itself provide a smaller-rank crossing term or a strict shortening.

\subsection{A genuine quadratic-weight kernel and its valuation cost}
The quadratic argument of the Racah and dual-Hahn polynomials is distinct from a quadratic grading in a generating parameter \cite[Table~18.25.1]{NISTDLMF}. Likewise, the Poisson kernel in \cite[equations~(1.6)--(1.10), pp.~1--2 of the preprint]{AlbaneseLawi2006} uses the linear index factor $e^{-nt}$, distinguished there from the heat eigenvalues. The Stieltjes--Wigert family does provide a bilateral quadratic weight \cite[equation~(2.2), Corollary~2.4 and Remark~2.5, pp.~2--5]{ChristiansenKoelink2006}. The following characteristic-two version is proved formally and does not import an analytic positivity assertion.

\begin{proposition}[A bilateral kernel and its monomial-gauge cost]\label{prop:bilateral-kernel}
Over $\mathbb F_2((t))$, let
\[
 P_n(x)=\sum_{j=0}^n t^{2j^2-j}\genfrac{[}{]}{0pt}{}{n}{j}_{t^2}x^j,
 \qquad
 B_{r,n}=\frac{t^nP_n(t^{2r})}{\prod_{j=1}^n(1+t^j)},
 \quad r\in\mathbb Z, n\geq0,
\]
where $\genfrac{[}{]}{0pt}{}{n}{j}_z$ is the Gaussian binomial polynomial. Then the bilateral sums converge coefficientwise and
\[
 \sum_{r\in\mathbb Z}t^{r^2}B_{r,i}B_{r,j}=\delta_{ij}.
\]
Their exact valuations are
\[
 \operatorname{val}_t B_{r,n}
 =n+\min_{0\leq j\leq n}\{2j^2+(2r-1)j\}.
\]
Suppose $c>0$ and $0\leq b<1/2$ satisfy
\[
 br^2+\tfrac c2n^2+\operatorname{val}_t B_{r,n}\geq0
 \quad(r\in\mathbb Z, n\geq0).
\]
Then $b\geq2/(c+4)$. In particular, with $a=1-2b>0$,
\[
 \frac ca\geq c+4>4.
\]
When the exponents lie in the chosen value group, the monomial gauge
$\widetilde B_{r,n}=t^{br^2+cn^2/2}B_{r,n}$
changes the Gram identity to
$\sum_r t^{ar^2}\widetilde B_{r,i}\widetilde B_{r,j}=\delta_{ij}t^{ci^2}$.
\end{proposition}
\begin{proof}
Define a functional on polynomials by $\mathcal L(x^m)=t^{-m^2}$. It has the bilateral representation
\[
 \sum_{r\in\mathbb Z}t^{r^2}(t^{2r})^m
 =t^{-m^2}\sum_{s\in\mathbb Z}t^{s^2}=t^{-m^2}.
\]
The terms indexed by $s$ and $-s$ cancel for $s\ne0$. Each Laurent coefficient receives only finitely many contributions. More generally, the same quadratic growth gives coefficientwise convergence after evaluating any fixed polynomial pair.

The finite Gaussian-binomial identity gives
\[
 \mathcal L(P_nx^k)
 =t^{-k^2}\prod_{j=0}^{n-1}(1+t^{2(j-k)}).
\]
This vanishes for $0\leq k<n$. Since the leading coefficient of $P_n$ is $t^{2n^2-n}$, taking $k=n$ shows
\[
 \mathcal L(P_n^2)=t^{-2n}\prod_{j=1}^n(1+t^{2j}).
\]
Division by its square root proves the stated orthogonality. All denominators are units at $t=0$.

Every Gaussian binomial polynomial has constant coefficient one. The difference between two candidate exponents is
\[
 \bigl(2j^2+(2r-1)j\bigr)
 -\bigl(2k^2+(2r-1)k\bigr)
 =(j-k)\bigl(2(j+k+r)-1\bigr),
\]
which is nonzero for distinct integers $j,k$. Hence the minimum is unique and cannot cancel, proving the valuation formula.

If $b=0$, take $n=1$ and let $r$ tend to negative infinity to contradict the required inequality. For $0<b<1/2$, choose integers $r_n$ with $r_n/n\to-1/b$. For sufficiently large $n$, the minimum in the valuation formula occurs at $j=n$. Divide the assumed inequality by $n^2$ and pass to the limit. This yields
$c/2+2-1/b\geq0$, or $b\geq2/(c+4)$. Consequently $a\leq c/(c+4)$, proving the ratio bound. The final Gram formula follows by substituting the monomial factors. The numerical inequality permits real $b,c$; actual formal series require the displayed exponents to belong to their value group, for example rational exponents after ramification.
\end{proof}

This kernel has genuine quadratic auxiliary weights. Its normalizer is one, so the original regular specialization with vanishing normalizer is absent. Moreover, within the stated monomial-gauge family, making all entries regular forces at least the displayed factor four between the target and auxiliary quadratic coefficients. A reciprocal crossing identity and a useful finite rank estimate for a different normalization remain unproved.
